\clearpage

\section{Phase Diagram}\label{sec-nbcp-phases}

The classical three-sublattice analysis supplies reference branches and
local boundaries. A phase diagram requires comparison with other
magnetic cells and the relevant quantum energies. The distinction is
maintained below: a stable magnon spectrum is a local test, while the
lowest energy among an incomplete set of candidate states need not
identify the global ground state.

\subsection{Classical phase diagram}\label{sec-nbcp-classical-phases}

\subsubsection{Three-sublattice classical
energy}\label{three-sublattice-classical-energy}

Let \(\mathbf n_\alpha=\mathbf S_\alpha/S\), with \(\alpha=A,B,C\),
specify a uniform three-sublattice configuration. Each pair of
sublattices is connected by all three bond orientations. Since
\(\sum_d\mathsf J_d=3\operatorname{diag}(J,J,J_z)\), its classical
energy per physical spin reduces to

\begin{equation}\phantomsection\label{eq-nbcp-three-energy}{
e_{\rm cl}=S^2\sum_{(\alpha\beta)=(AB),(BC),(CA)}
\left[J\sin\vartheta_\alpha\sin\vartheta_\beta
\cos(\phi_\alpha-\phi_\beta)
+J_z\cos\vartheta_\alpha\cos\vartheta_\beta\right]
-\frac{hS}{3}\sum_\alpha\cos\vartheta_\alpha.
}\end{equation}

Here \(\vartheta\) may be signed, with
\(\mathbf n=(\sin\vartheta\cos\phi,\sin\vartheta\sin\phi,\cos\vartheta)\).
Negative \(\vartheta\) is a coordinate choice, equivalent to a positive
polar angle and an azimuth shifted by \(\pi\).

The SOC contributions cancel in this restricted classical energy. They
do not disappear from the operator Hamiltonian, from finite-momentum
fluctuations, or from states with other magnetic unit cells. In
particular, the cancellation does not establish SOC-independent global
phase boundaries.

\subsubsection{Classical reference fields and their stability
meaning}\label{classical-reference-fields-and-their-stability-meaning}

The UUD stability matrix can be derived in regular Cartesian coordinates
at the collinear poles. Put
\(\mathbf S_A=(\sqrt S x_A,\sqrt S y_A,S-(x_A^2+y_A^2)/2)\) and
similarly for B; for C replace the last component by
\(-S+(x_C^2+y_C^2)/2\). To quadratic order,
\(\delta e=(\mathbf x^{\mathsf T}M\mathbf x+\mathbf y^{\mathsf T}M\mathbf y)/2\),
where

\begin{equation}\phantomsection\label{eq-nbcp-uud-hessian}{
M=\begin{pmatrix}
h/3&JS&JS\\
JS&h/3&JS\\
JS&JS&2J_zS-h/3
\end{pmatrix}.
}\end{equation}

The antisymmetric A/B eigenvector has eigenvalue \(h/3-JS\). The
symmetric A/B sector and C reduce to a \(2\times2\) matrix with diagonal
entries \(h/3+JS\) and \(2J_zS-h/3\), and off-diagonal entry
\(\sqrt2JS\). Setting the lower eigenvalue to zero gives the upper UUD
boundary. At P, the diagonal entries become \(h/3-2J_zS\), giving the
saturation boundary. For \(J_z>J>0\) the three reference fields are

\begin{equation}\phantomsection\label{eq-nbcp-critical-fields}{
\begin{aligned}
h_{c1}&=3SJ,\\
h_{c2}&=3S\left[J_z-\frac J2+
\sqrt{J_z^2+J_zJ-\frac74J^2}\right],\\
h_{c3}&=3S(J+2J_z),\qquad B_{ci}=h_{ci}/(g_z\mu_B).
\end{aligned}
}\end{equation}

These are the classical XXZ reference boundaries and uniform
three-sublattice softening conditions. Quantum corrections, other
ordering wave vectors and SOC-driven competitors require separate
calculations.

\begin{longtable}[]{@{}
  >{\raggedleft\arraybackslash}p{(\linewidth - 4\tabcolsep) * \real{0.3333}}
  >{\raggedright\arraybackslash}p{(\linewidth - 4\tabcolsep) * \real{0.3333}}
  >{\raggedright\arraybackslash}p{(\linewidth - 4\tabcolsep) * \real{0.3333}}@{}}
\toprule\noalign{}
\begin{minipage}[b]{\linewidth}\raggedleft
\(J\) (meV)
\end{minipage} & \begin{minipage}[b]{\linewidth}\raggedright
\((h_{c1},h_{c2},h_{c3})\) (meV)
\end{minipage} & \begin{minipage}[b]{\linewidth}\raggedright
\((B_{c1},B_{c2},B_{c3})\) (T)
\end{minipage} \\
\midrule\noalign{}
\endfirsthead
\toprule\noalign{}
\begin{minipage}[b]{\linewidth}\raggedleft
\(J\) (meV)
\end{minipage} & \begin{minipage}[b]{\linewidth}\raggedright
\((h_{c1},h_{c2},h_{c3})\) (meV)
\end{minipage} & \begin{minipage}[b]{\linewidth}\raggedright
\((B_{c1},B_{c2},B_{c3})\) (T)
\end{minipage} \\
\midrule\noalign{}
\endhead
\bottomrule\noalign{}
\caption{Classical three-sublattice reference fields for the stated exchange parameters; the starred row is the literature set.}\label{tab-nbcp-03-phase-diagram-1}\\
\endlastfoot
0.075 & (0.112500, 0.315916, 0.487500) & (0.418417, 1.174976,
1.813142) \\
0.076 & (0.114000, 0.314316, 0.489000) & (0.423996, 1.169024,
1.818721) \\
0.0779\(^{\ast}\) & (0.116850, 0.302359, 0.484350) & (0.428053, 1.107621,
1.774305) \\
\end{longtable}

The first two rows use \(J_z=0.125\) meV and \(g_z=4.645\). The starred
row is the literature set of Woodland et al.~\cite{woodland2025}, Table~1,
with \(J_z=0.1225\) meV and \(g_c=4.716\); there \(J\) is their
\(J_{xy}\), the coefficient of \(\hat S^x\hat S^x+\hat S^y\hat S^y\).
All rows use \(S=1/2\) and
\(\mu_B=0.05788381806\) meV/T. The older notes also quote
\((0.1005,0.3278,0.4755)\) meV while specifying \(J=0.076\) meV. Those
numbers are not evaluations of Eq.~\ref{eq-nbcp-critical-fields} at the
stated inputs and are not used here. The present reference fields were
recalculated; the archived phase scans were not rerun.

\subsection{Quantum corrections: zero-point
energy}\label{sec-nbcp-phase-quantum-correction}

At harmonic order, the comparison adds the spin-wave vacuum energy to
the classical energy for each admissible reference state. The positive
physical magnon bands and the normal-ordering trace constant must both
be included, and all candidate energies must use the same per-spin
normalization and Brillouin-zone convention. An unstable quadratic
background cannot be assigned a physical harmonic vacuum energy by
discarding its unstable modes.

The existing calculations resolve the vacuum-energy dependence along
selected Y/V angular orbits; they do not reoptimize and compare all
candidate phases across the model's parameter space. Consequently, the
classical reference fields are not promoted to quantum-corrected phase
boundaries. The local angular selection and its pseudo-Goldstone gap are
treated in Section~\nbcpnoteref{Y/V low energy and pseudo-Goldstone gap}{``Y and V Phases: Symmetry and Low-Energy Theory''}; numerical quadrature
and convergence are recorded in Section~\nbcpnoteref{Numerical verification and provenance}{``Numerical convergence''}.
